\section{The trees of the child graph}
\label{sec:trees}

Section~11 of~\cite{comma} reports that, of the 50 trees of the base-10 child graph, only the tree
rooted at $20$ is infinite, and that the longest-lived rival, rooted at $30$, persists until
$10^{365} - 82$. As a first check of our implementation we reproduced this by walking every tree
completely, using the child-graph jump bound of Section~\ref{sec:jumps} for unbranched stretches.
Since paths never merge, each node is visited at most once, and the cost grows with the number of
digits reached and the number of branch points rather than with the number of nodes.

\begin{proposition}[{cf.~\cite[\S11]{comma}}]
  Following all paths up to $1000$ digits, every tree of the base-10 child graph except the one
  rooted at $20$ is finite. The tree rooted at $30$ has $1007$ branch points and $1008$ leaves; its
  largest leaf is exactly $10^{365} - 82$, reached after about $2 \cdot 10^{363}$ steps from the
  root. No other tree reaches $10^{133}$.
\end{proposition}

\begin{remark}
  Section~11 of~\cite{comma} says that the root-$30$ path ``persisted for $10^{365} - 82$ terms''.
  The number $10^{365} - 82$ is the value of the landmine; the path itself has about
  $2 \cdot 10^{363}$ terms, in line with the rule of thumb that the $n$-th term is about $50 n$.
\end{remark}

Table~\ref{tab:trees} lists the trees with at least ten leaves. All 49 finite trees are listed in
Appendix~\ref{app:trees}. Trees without branch points are single successor paths, and for the roots
$3, 4, 7$ they reproduce the corresponding entries of (1.2) and (1.3).

\begin{table}[ht]
  \centering
  \begin{tabular}{@{}rrrcc@{}}
    \toprule
    root & leaves & branch points & height (edges)                   & largest leaf              \\
    \midrule
       6 &     32 &            31 & $\approx 10^{40.31\phantom{0}}$  & $10^{42\phantom{0}} - 64$ \\
      13 &     66 &            65 & $\approx 10^{85.31\phantom{0}}$  & $10^{87\phantom{0}} - 37$ \\
      17 &     31 &            30 & $\approx 10^{52.32\phantom{0}}$  & $10^{54\phantom{0}} - 55$ \\
      30 &   1008 &          1007 & $\approx 10^{363.30}$            & $10^{365}           - 82$ \\
      38 &     10 &             9 & $\approx 10^{15.32\phantom{0}}$  & $10^{17\phantom{0}} - 37$ \\
      64 &    285 &           284 & $\approx 10^{130.33}$            & $10^{132}           - 19$ \\
    \bottomrule
  \end{tabular}
  \caption{The finite trees of the base-10 child graph with at least ten leaves.}
  \label{tab:trees}
\end{table}

Every positive integer lies in exactly one tree. As the 49 finite trees contain no number above
$10^{365}$, the tree rooted at $20$ contains \emph{every} number above $10^{365}$. In particular,
above this bound its successor paths are exactly the successor sequences starting at the numbers
$c \cdot 10^j$, which motivates the next section.
